% Introduction to the king and pawn chapter, for the 1500-1599 book only.
% The theory it needs comes from the library (lib/), placed by make_book where it is first used.

\section*{What to look for}
Pawn endings are the simplest endings, and every other ending can turn into one. When players at
this level go wrong in them, it is rarely for lack of theory: they move too fast. Three habits would have saved most of the games
in this chapter.

Count before you race. When both sides run with their pawns, count the moves: how many does
my pawn need to queen, how many does his, and who moves first? Does my new queen come with check?
Counting\index{Practical play!Counting} takes a minute and decides the game, as in
\xref{Diagram}{yougaram-riveranegron-60w}.

Do not weaken yourself. A pawn move cannot be taken back. If your position holds as it is,
a waiting king move is often better than a pawn move that gives the enemy king a new
square,\index{Strategy!Do not weaken yourself} as in \xref{Diagram}{santiagollanes-amadorgarcia-57w}.

Rule out what loses. In many positions you do not need to find the best move; it is enough to
see that the other moves lose. The process of elimination\index{Practical play!Process of elimination}
and simple logic\index{Practical play!Simple logic} find the right move without long calculation
(\xref{Diagram}{ince-tokat-80w} and \xref{Diagram}{ince-tokat-87b}). And before you make your plan,
ask what your opponent wants: a little prophylaxis\index{Strategy!Prophylaxis} often decides the race
(\xref{Diagram}{perezselios-fuentes-58w}).

Watch out for stalemate, too: a king in the corner in front of a rook's pawn holds many positions
that look lost.

